\documentclass[12pt,tikz,border=4mm]{standalone}

\usepackage{fontspec}
\usepackage{xeCJK}
\usepackage{xcolor}
\usetikzlibrary{arrows.meta,backgrounds,calc,fit,positioning}

\setsansfont{Noto Sans CJK SC}
\setCJKsansfont{Noto Sans CJK SC}
\renewcommand{\familydefault}{\sfdefault}

\definecolor{signal}{HTML}{009E73}
\definecolor{prompt}{HTML}{0072B2}
\definecolor{delayed}{HTML}{D55E00}
\definecolor{shared}{HTML}{4E5D6C}
\definecolor{ink}{HTML}{1F2933}

\begin{document}
\begin{tikzpicture}[
  x=1cm,
  y=1cm,
  text=ink,
  font=\fontsize{11.8pt}{13.6pt}\selectfont,
  flow/.style={-{Latex[length=2.5mm,width=1.7mm]}, draw=ink!72,
               line width=0.85pt, rounded corners=2pt},
  dep/.style={-{Latex[length=2.2mm,width=1.5mm]}, draw=shared!72,
              line width=0.75pt, densely dashed, rounded corners=2pt},
  signalbox/.style={rectangle, rounded corners=3pt, draw=signal,
                   fill=signal!8, line width=0.9pt, align=center,
                   text width=2.55cm, minimum height=1.08cm, inner sep=4pt},
  promptbox/.style={rectangle, rounded corners=3pt, draw=prompt,
                   fill=prompt!8, line width=0.9pt, align=center,
                   text width=2.95cm, minimum height=1.08cm, inner sep=4pt},
  delayedbox/.style={rectangle, rounded corners=3pt, draw=delayed,
                    fill=delayed!8, line width=0.9pt, align=center,
                    text width=2.75cm, minimum height=1.08cm, inner sep=4pt},
  modelbox/.style={rectangle, rounded corners=3pt, draw=shared!80,
                  fill=shared!8, line width=0.85pt, align=center,
                  text width=2.7cm, minimum height=0.9cm, inner sep=3pt},
  eventbox/.style={rectangle, rounded corners=3pt, draw=#1,
                  fill=white, densely dashed, line width=0.9pt, align=center,
                  text width=2.25cm, minimum height=0.92cm, inner sep=3pt},
  sharedbox/.style={rectangle, rounded corners=4pt, draw=shared,
                   fill=shared!7, line width=1.0pt, align=center,
                   text width=4.65cm, minimum height=1.08cm, inner sep=5pt},
  outputbox/.style={rectangle, rounded corners=3pt, draw=shared!85,
                   fill=white, line width=0.9pt, align=center,
                   text width=2.35cm, minimum height=1.0cm, inner sep=4pt},
  lanehead/.style={rounded corners=3pt, text=white, align=center,
                  font=\fontsize{12.2pt}{14pt}\selectfont\bfseries,
                  minimum width=1.55cm, minimum height=1.0cm, inner sep=4pt},
  smallnote/.style={font=\fontsize{9.7pt}{11.3pt}\selectfont,
                   text=ink!82, align=left}
]

% Title and shared detector-model contract
\node[font=\fontsize{16pt}{18pt}\selectfont\bfseries, text=ink]
  at (12.75,1.1) {端到端模拟与分析流程（图 2 候选稿）};

\node[modelbox, text width=15.4cm, minimum height=0.72cm] (detmodel)
  at (12.75,0.15)
  {共享探测器模型：探测器--低温系统质量模型用于信号重放、瞬发输运、活化产生和延迟衰变输运};

% Lane backgrounds
\begin{scope}[on background layer]
  \filldraw[rounded corners=5pt, fill=signal!3.5, draw=signal!52, line width=0.8pt]
    (0.15,-0.55) rectangle (17.1,-3.15);
  \filldraw[rounded corners=5pt, fill=prompt!3.5, draw=prompt!52, line width=0.8pt]
    (0.15,-3.45) rectangle (17.1,-7.25);
  \filldraw[rounded corners=5pt, fill=delayed!3.5, draw=delayed!52, line width=0.8pt]
    (0.15,-7.55) rectangle (17.1,-13.05);
  \filldraw[rounded corners=5pt, fill=shared!3.5, draw=shared!52, line width=0.8pt]
    (17.45,-0.55) rectangle (25.35,-13.05);
\end{scope}

% Lane labels
\node[lanehead, fill=signal] at (1.12,-1.85) {聚焦信号};
\node[lanehead, fill=prompt, minimum height=1.35cm] at (1.12,-5.35)
  {瞬发外部\\大气本底};
\node[lanehead, fill=delayed, minimum height=1.35cm] at (1.12,-10.3)
  {延迟活化\\本底};
\node[lanehead, fill=shared, minimum height=1.35cm] at (18.35,-1.45)
  {统一处理\\与输出};

% Focused-signal lane
\node[signalbox, text width=2.15cm] (src) at (3.05,-1.25)
  {511 keV\\参考点源};
\node[modelbox, text width=2.15cm] (optmodel) at (3.05,-2.45)
  {Laue 光学\\质量模型};
\node[signalbox] (optics) at (6.15,-1.85) {Laue 光学输运};
\node[signalbox] (phase) at (9.25,-1.85) {焦平面聚焦光子\\相空间};
\node[signalbox] (replay) at (12.35,-1.85) {探测器耦合\\信号重放};
\node[eventbox=signal] (sigevt) at (15.35,-1.85) {聚焦信号\\事例目录};

\draw[flow] (src.east) -| ($(optics.west)+(0,0.28)$);
\draw[dep] (optmodel.east) -| ($(optics.west)+(0,-0.28)$);
\draw[flow] (optics) -- (phase);
\draw[flow] (phase) -- (replay);
\draw[flow] (replay) -- (sigevt);

% Prompt-background lane: two computational streams inside one physical class
\node[promptbox, text width=3.05cm] (expacs) at (3.55,-4.45)
  {EXPACS/PARMA\\大气 $\gamma$ 连续谱与\\瞬发粒子场};
\node[promptbox] (prompttr) at (7.35,-4.45)
  {PROMPT\\瞬发直接输运};
\node[eventbox=prompt] (promptevt) at (15.35,-4.45)
  {宽带瞬发\\事例目录};

\node[promptbox, text width=3.05cm] (linesrc) at (3.55,-6.2)
  {半经验大气正电子湮没\\511 keV 线源};
\node[promptbox] (linetr) at (7.35,-6.2)
  {独立单能\\探测器输运};
\node[eventbox=prompt] (lineevt) at (15.35,-6.2)
  {大气 511 keV\\瞬发事例目录};

\draw[flow] (expacs) -- (prompttr);
\draw[flow] (prompttr) -- (promptevt);
\draw[flow] (linesrc) -- (linetr);
\draw[flow] (linetr) -- (lineevt);

% Delayed-activation lane: independent BUILDUP chain
\node[delayedbox, text width=3.05cm] (expacs2) at (3.55,-8.75)
  {EXPACS/PARMA\\大气粒子场\\（与瞬发分支相同）};
\node[delayedbox] (buildup) at (7.35,-8.75)
  {BUILDUP\\活化产生输运};
\node[delayedbox, text width=3.05cm] (records) at (11.15,-8.75)
  {核素产额、材料体积\\与产生位置记录};

\node[delayedbox, text width=3.05cm] (inventory) at (9.75,-10.65)
  {飞行时间活度积累与\\参考时刻核素清单};
\node[delayedbox, text width=2.55cm] (dsrc) at (13.0,-10.65)
  {按产生位置构建\\延迟衰变源};
\node[delayedbox, text width=2.25cm] (decaytr) at (13.0,-12.15)
  {延迟衰变输运};
\node[eventbox=delayed, text width=2.1cm] (delayevt) at (15.55,-12.15)
  {延迟活化\\事例目录};

\draw[flow] (expacs2) -- (buildup);
\draw[flow] (buildup) -- (records);
\draw[flow] (records) -- (inventory);
\draw[flow] (inventory) -- (dsrc);
\draw[flow] (dsrc) -- (decaytr);
\draw[flow] (decaytr) -- (delayevt);

% Event-stream bus and common offline response/selection
\coordinate (busTop) at (17.2,-1.85);
\coordinate (busMid) at (17.2,-5.35);
\coordinate (busBottom) at (17.2,-12.15);
\draw[draw=shared!72, line width=1.15pt] (busTop) -- (busBottom);
\draw[flow] (sigevt.east) -- (busTop);
\draw[flow] (promptevt.east) -- (17.2,-4.45);
\draw[flow] (lineevt.east) -- (17.2,-6.2);
\draw[flow] (delayevt.east) -- (busBottom);

\node[sharedbox, text width=5.25cm, minimum height=3.55cm] (selection)
  at (21.45,-4.95)
  {\textbf{统一探测器响应与事例选择}\\[2pt]
   按 TES 像素汇总能量沉积\\
   420 eV FWHM 响应与 0.3 keV 读出阈值\\
   $W_{511}=510.58$--$511.42$ keV\\
   BGO 总沉积 $\geq 50$ keV 的离线主动反符合\\
   康普顿拓扑/视场一致性检验};
\draw[flow] (busMid) -- (selection.west);

\node[sharedbox, text width=4.55cm] (rates) at (21.45,-8.15)
  {选后信号率与分量本底率};
\draw[flow] (selection) -- (rates);

% Two analysis/output branches
\node[outputbox, text width=2.65cm] (decomp) at (19.35,-10.15)
  {选后本底来源分解\\[-1pt]
   \fontsize{9.5pt}{11pt}\selectfont 入射粒子族/母核素/产生体积};
\node[outputbox, text width=2.65cm] (shield) at (19.35,-12.15)
  {屏蔽几何设计与比较};

\node[outputbox, text width=2.75cm] (fold) at (23.45,-10.15)
  {任务时间--轨迹\\解析率折叠};
\node[outputbox, text width=2.75cm] (performance) at (23.45,-12.15)
  {累计信号、本底计数\\与未分辨线灵敏度};

\draw[flow] (rates.south) -- ++(0,-0.42) -| (decomp.north);
\draw[flow] (rates.south) -- ++(0,-0.42) -| (fold.north);
\draw[flow] (decomp) -- (shield);
\draw[flow] (fold) -- (performance);

% Figure-level notes
\node[smallnote, text width=24.6cm, draw=shared!38, fill=white,
      rounded corners=2pt, inner sep=5pt] at (12.75,-13.85)
  {\textbf{说明：}大气 511 keV 线在物理分类上属于瞬发外部大气光子本底，但采用独立的半经验源模型、单能输运和任务标度。图中的 50 keV 是后处理中对主动屏蔽体积逐事例汇总能量沉积所施加的离线反符合阈值，而不是原生或硬件触发阈值。};

\end{tikzpicture}
\end{document}
